\documentclass[12pt,oneside]{book}

% =====================================================
% METADATOS
% =====================================================
\newcommand{\universidad}{Universidad de Buenos Aires (UBA)}
\newcommand{\materia}{Derecho Procesal Civil y Comercial}
\newcommand{\titulo}{Contestación de Demanda, Actitudes del Demandado, Rebeldía, Allanamiento y Reconvención}
\newcommand{\autor}{Isaac Edgar Camacho Ocampo}
\newcommand{\anio}{2020}

% =====================================================
% PAQUETES
% =====================================================
\usepackage[utf8]{inputenc}
\usepackage[T1]{fontenc}
\usepackage[spanish,es-nodecimaldot]{babel}

\usepackage[
    top=2.5cm,
    bottom=2.5cm,
    left=3cm,
    right=2.5cm,
    headheight=15pt
]{geometry}

\usepackage{amsmath}
\usepackage{amssymb}
\usepackage{mathtools}

\usepackage{graphicx}
\usepackage{float}
\usepackage{caption}
\usepackage{subcaption}

\usepackage{booktabs}
\usepackage{array}
\usepackage{longtable}
\usepackage{tabularx}

\usepackage{enumitem}

\usepackage{xcolor}
\usepackage[most]{tcolorbox}

\usepackage{fancyhdr}
\usepackage{ragged2e}

\usepackage[
    colorlinks=true,
    linkcolor=blue!50!black,
    citecolor=green!40!black,
    urlcolor=blue!60!black,
    pdftitle={\titulo},
    pdfauthor={\autor},
    pdfsubject={Apuntes universitarios},
    bookmarks=true
]{hyperref}

% =====================================================
% RUTAS DE IMÁGENES
% =====================================================
\graphicspath{
    {./img/}
    {./graficos/}
    {./figuras/}
}

% =====================================================
% CONFIGURACIÓN GENERAL
% =====================================================
\setcounter{tocdepth}{2}

\setlength{\parindent}{0pt}
\setlength{\parskip}{0.5em}

\setlist[itemize]{
    topsep=0.3em,
    itemsep=0.2em
}

\setlist[enumerate]{
    topsep=0.3em,
    itemsep=0.2em
}

\clubpenalty=10000
\widowpenalty=10000

% =====================================================
% COMANDOS MATEMÁTICOS
% =====================================================
\newcommand{\abs}[1]{\left|#1\right|}
\newcommand{\norm}[1]{\left\lVert#1\right\rVert}
\newcommand{\vect}[1]{\mathbf{#1}}
\newcommand{\deriv}[2]{\frac{d#1}{d#2}}
\newcommand{\pderiv}[2]{\frac{\partial#1}{\partial#2}}

% =====================================================
% ENTORNOS / CAJAS ACADÉMICAS
% =====================================================
\newtcolorbox{definition}[1]{
    enhanced,
    breakable,
    title={Definición: #1},
    fonttitle=\bfseries,
    colback=blue!3,
    colframe=blue!50!black,
    boxrule=0.8pt,
    arc=2mm
}

\newtcolorbox{important}{
    enhanced,
    breakable,
    title={Importante},
    fonttitle=\bfseries,
    colback=yellow!8,
    colframe=orange!70!black,
    boxrule=0.8pt,
    arc=2mm
}

\newtcolorbox{example}[1]{
    enhanced,
    breakable,
    title={Ejemplo: #1},
    fonttitle=\bfseries,
    colback=green!4,
    colframe=green!40!black,
    boxrule=0.8pt,
    arc=2mm
}

\newtcolorbox{formula}[1]{
    enhanced,
    breakable,
    title={Fórmula / Regla: #1},
    fonttitle=\bfseries,
    colback=purple!4,
    colframe=purple!50!black,
    boxrule=0.8pt,
    arc=2mm
}

\newtcolorbox{exercise}[1]{
    enhanced,
    breakable,
    title={Ejercicio: #1},
    fonttitle=\bfseries,
    colback=gray!5,
    colframe=gray!60!black,
    boxrule=0.8pt,
    arc=2mm
}

\newtcolorbox{note}{
    enhanced,
    breakable,
    title={Nota},
    fonttitle=\bfseries,
    colback=white,
    colframe=gray!50,
    boxrule=0.6pt,
    arc=2mm
}

% =====================================================
% ENCABEZADOS Y PIES
% =====================================================
\pagestyle{fancy}
\fancyhf{}

\fancyhead[L]{\nouppercase{\leftmark}}
\fancyhead[R]{\materia}

\fancyfoot[C]{\thepage}

\renewcommand{\headrulewidth}{0.4pt}
\renewcommand{\footrulewidth}{0pt}

\fancypagestyle{plain}{
    \fancyhf{}
    \fancyfoot[C]{\thepage}
    \renewcommand{\headrulewidth}{0pt}
}

% =====================================================
% DOCUMENTO
% =====================================================
\begin{document}

% -----------------------------------------------------
% PORTADA
% -----------------------------------------------------
\begin{titlepage}

\thispagestyle{empty}

\begin{center}

\vspace*{1.5cm}

{\large \textsc{\universidad}}

\vspace{3cm}

{\Huge \textbf{\materia}}

\vspace{1cm}

{\LARGE \titulo}

\vspace{0.6cm}

{\normalsize Unidad 10 — Cátedra a cargo del Dr. Jorge Rojas}

\vspace{0.8cm}

\textit{Comprender los conceptos es el primer paso para convertir el estudio en conocimiento.}

\vspace{1.2cm}

\rule{0.70\textwidth}{0.5pt}

\vspace{0.5cm}

{\large Apuntes de estudio}

\vspace{0.5cm}

\rule{0.70\textwidth}{0.5pt}

\vfill

{\large \autor}

\vspace{0.3cm}

{\large \anio}

\end{center}

\vfill

\begin{flushleft}
\textit{Este es un modesto aporte a los estudiantes de ``Derecho Procesal Civil y Comercial'' no pretende sustituir el material oficial de la materia.}
\end{flushleft}

\end{titlepage}

% -----------------------------------------------------
% ÍNDICE
% -----------------------------------------------------
\tableofcontents

% =====================================================
% INTRODUCCIÓN A LA UNIDAD
% =====================================================
\chapter*{Presentación de la unidad}
\addcontentsline{toc}{chapter}{Presentación de la unidad}

Así como en un juicio el proceso avanza etapa por etapa sin poder retroceder, en virtud del principio de preclusión, el desarrollo de esta unidad sigue también un orden lógico: en primer lugar se repasa la ubicación de la Unidad 10 dentro del proceso judicial; luego se estudia en detalle cuáles son las actitudes que puede asumir el demandado al recibir la demanda; y finalmente se analizan las consecuencias de cada una de esas decisiones —contestar, allanarse, reconvenir o quedar en rebeldía—.

\textbf{Temas tratados en esta unidad:}
\begin{itemize}
    \item Repaso de las etapas del proceso judicial y ubicación de la Unidad 10.
    \item La contestación de demanda: concepto, requisitos formales y sustanciales.
    \item Artículo 356 CPCCN: cargas del demandado al contestar.
    \item Instrumentos públicos y privados: diferencias en su impugnación.
    \item Actitudes omisivas y positivas del demandado.
    \item Documentos nuevos (Art. 335) y hechos no invocados (Art. 334).
    \item Allanamiento: concepto, tipos y oportunidad.
    \item Reconvención: concepto, requisitos, efectos y trámite.
    \item Rebeldía: declaración, efectos, medidas cautelares y cesación.
\end{itemize}

% =====================================================
\chapter{Repaso del proceso judicial y contexto}
% =====================================================

\section{Contexto de la unidad}

La presente unidad se ubica luego del estudio de los sujetos procesales y del acto procesal de la demanda, y tiene por objeto el estudio de la contestación de demanda y de las restantes actitudes que puede asumir el demandado frente a ella.

\section{Repaso: las etapas del proceso judicial}

El legislador optó por un sistema de preclusión, y no de unicidad, en virtud del cual el proceso se diagrama en distintas etapas sucesivas.

\begin{important}
\textbf{Principio de preclusión:} las partes pueden ejercer determinados actos en cada etapa del proceso, en determinada forma y en determinado tiempo. Si no los realizan en esa etapa o en ese tiempo, pierden la posibilidad de hacerlos más adelante. El proceso avanza siempre hacia el dictado de la sentencia y nunca se retrotrae, salvo que exista un planteo de nulidad.
\end{important}

Las etapas del proceso judicial son las siguientes:

\begin{table}[H]
\centering
\caption{Etapas del proceso judicial}
\begin{tabularx}{\textwidth}{l X}
\toprule
\textbf{Etapa} & \textbf{Descripción} \\
\midrule
Introductoria & Se inicia el proceso: demanda, traslado, contestación, excepciones. \\
Probatoria & Las partes ofrecen y producen la prueba ofrecida en los escritos introductorios. \\
Conclusional & Las partes presentan alegatos y el juez analiza las pruebas. \\
Sentencia & El juez dicta sentencia: resuelve el conflicto y actúa la voluntad de la ley. \\
Recursiva & Las partes pueden impugnar la sentencia mediante el recurso de apelación ante la Cámara. \\
De ejecución & Si el condenado no cumple voluntariamente con la sentencia, se procede a la ejecución forzada. \\
\bottomrule
\end{tabularx}
\end{table}

\section{Los actos procesales y su clasificación}

Los actos procesales deben cumplir ciertas formas y realizarse en determinado tiempo, y se clasifican en tres grandes categorías:

\begin{table}[H]
\centering
\caption{Clasificación de los actos procesales}
\begin{tabularx}{\textwidth}{l X X}
\toprule
\textbf{Tipo de acto} & \textbf{Definición} & \textbf{Ejemplo} \\
\midrule
De postulación & Vinculados con el ejercicio del derecho de postulación de las partes. Son los de mayor trascendencia en el proceso, ya que propenden hacia el dictado de la sentencia. & La demanda y la contestación de demanda. \\
De obtención & Hacen al mero trámite del proceso. Las partes solicitan al juez determinadas resoluciones para avanzar en el proceso. & Solicitar al juez que libre una nueva cédula con habilitación de días y horas inhábiles para notificar al demandado. \\
De comunicación & Tienen que ver con las notificaciones entre los sujetos procesales. & La cédula de notificación del traslado de la demanda. \\
\bottomrule
\end{tabularx}
\end{table}

\section{Sistemas alternativos de resolución de conflictos (repaso)}

Antes de iniciar el proceso, existen sistemas alternativos de resolución de conflictos:

\begin{table}[H]
\centering
\caption{Sistemas alternativos de resolución de conflictos}
\begin{tabularx}{\textwidth}{l X}
\toprule
\textbf{Sistema} & \textbf{Características} \\
\midrule
Negociación directa & Las partes pueden arribar de modo directo a un acuerdo que ponga fin al conflicto, sin necesidad de iniciar el proceso judicial. \\
Mediación (previa y obligatoria) & Es un método alternativo de resolución de conflictos que es previo y obligatorio al proceso judicial. Quien quiera promover una demanda debe cumplir primero con esta instancia, convocando a quien pretende demandar para intentar arribar a un acuerdo. Si no se llega a un acuerdo, se labra un acta y queda habilitada la instancia judicial. \\
\bottomrule
\end{tabularx}
\end{table}

\section{La demanda (repaso de la unidad anterior)}

\begin{definition}{La demanda}
Es un acto procesal de postulación en virtud del cual la parte actora ejerce una pretensión hacia el demandado para que éste la satisfaga. Contiene una pretensión y tiene por efecto iniciar el proceso judicial, con la finalidad de obtener una sentencia en la cual se haga actuar la voluntad de la ley y se reconozca el derecho del actor.
\end{definition}

\begin{example}{Tipos de pretensiones en una demanda}
\begin{enumerate}
    \item Proceso de desalojo: la pretensión del actor es que el demandado deje el inmueble de su propiedad.
    \item Daños y perjuicios por accidente de tránsito: la pretensión es obtener una indemnización del responsable del accidente.
    \item Daños y perjuicios por mala praxis médica: la pretensión es obtener una indemnización por los daños ocasionados en una cirugía.
\end{enumerate}
\end{example}

\section{Actitudes del demandado tras la notificación: introducción}

Una vez que el actor promueve la demanda y cumple ese acto procesal con todos los requisitos legales, el juzgado dicta una resolución en la que tiene al actor por parte, por presentado y por constituido el domicilio, y ordena correr traslado de esa demanda al demandado para que éste se defienda.

La demanda se notifica generalmente por cédula dirigida al domicilio real del demandado. También puede notificarse por edictos, por radiodifusión, personalmente o por acta notarial, aunque la cédula es el medio habitual.

La cédula de notificación informa al demandado los siguientes datos:

\begin{table}[H]
\centering
\caption{Datos que informa la cédula de notificación}
\begin{tabularx}{\textwidth}{l X}
\toprule
\textbf{Dato} & \textbf{Descripción} \\
\midrule
Juzgado interviniente & El juzgado que está a cargo del proceso. \\
Carátula del expediente & El nombre del expediente (ej.: ``Gómez, María Eugenia c/ Pérez, Juan s/ Daños y Perjuicios''). \\
Número de expediente & Cada expediente tiene un número y el año (ej.: 2020/2020). \\
\bottomrule
\end{tabularx}
\end{table}

\section{Las cargas que nacen con la notificación}

Con la notificación nacen en cabeza del demandado dos cargas:

\begin{table}[H]
\centering
\caption{Cargas procesales del demandado}
\begin{tabularx}{\textwidth}{l X}
\toprule
\textbf{Carga} & \textbf{Descripción} \\
\midrule
Comparecer & Presentarse ante la jurisdicción en respuesta al llamado (citación). \\
Contestar & Contestar la demanda dentro del plazo legal (emplazamiento). \\
\bottomrule
\end{tabularx}
\end{table}

\begin{important}
\textbf{Carga procesal vs. deber procesal:} la diferencia es que la carga es ``el imperativo del propio interés'': si la parte no la ejerce, no hay sanción, pero el no ejercicio puede beneficiar al adversario en el proceso. El deber, en cambio, sí trae aparejada una sanción.
\end{important}

\section{Clasificación de las actitudes del demandado}

\begin{table}[H]
\centering
\caption{Actitudes del demandado frente a la demanda}
\begin{tabularx}{\textwidth}{l l X}
\toprule
\textbf{Tipo} & \textbf{Actitud} & \textbf{Descripción} \\
\midrule
Omisiva & No comparece y no contesta & El demandado no cumple con ninguna de las dos cargas. Conduce a la rebeldía (Arts. 59 y ss. CPCCN). \\
Omisiva & Comparece pero no contesta & El demandado cumple sólo con la carga de comparecer, presentando un escrito donde denuncia su domicilio real, constituye domicilio procesal y electrónico, y nada más. No es lo mismo que la rebeldía: puede ejercer todos los actos procesales de ahí en adelante. \\
Positiva & Contesta (defensiva) & Cumple ambas cargas. Pide que se rechace la pretensión del actor (Art. 356 CPCCN). \\
Positiva & Contesta y opone excepciones & Además de contestar, opone las defensas de los Arts. 346, 347 y 348 CPCCN. \\
Positiva & Contesta y se allana & Se somete a la pretensión del actor, total o parcialmente. \\
Positiva & Contesta y reconviene & Además de contestar, ejerce una pretensión propia contra el actor dentro del mismo proceso. \\
\bottomrule
\end{tabularx}
\end{table}

\section*{Cuadro Cornell del capítulo}
\addcontentsline{toc}{section}{Cuadro Cornell del capítulo}

\begin{table}[H]
\centering
\begin{tabularx}{\textwidth}{
    >{\raggedright\arraybackslash}p{0.30\textwidth}
    >{\raggedright\arraybackslash}X
}
\toprule
\textbf{Preguntas / palabras clave} & \textbf{Notas / respuestas} \\
\midrule
\textbf{¿Qué sistema adoptó el legislador para organizar el proceso y qué implica?} & Adoptó el sistema de preclusión: el proceso se diagrama en etapas sucesivas y cada acto debe realizarse en su oportunidad; de lo contrario se pierde la posibilidad de realizarlo después. \\
\textbf{¿Cuáles son las etapas del proceso judicial?} & Introductoria, probatoria, conclusional, sentencia, recursiva y de ejecución. \\
\textbf{¿Cuáles son las dos cargas que nacen con la notificación de la demanda?} & Comparecer (presentarse ante la jurisdicción) y contestar (responder la demanda dentro del plazo legal). \\
\textbf{¿Qué diferencia hay entre carga y deber procesal?} & La carga es el imperativo del propio interés: su incumplimiento no trae sanción, pero puede beneficiar al adversario. El deber, en cambio, sí trae aparejada una sanción. \\
\midrule
\multicolumn{2}{p{\textwidth}}{
\textbf{Resumen:} El proceso judicial se organiza en etapas sucesivas bajo el principio de preclusión. Notificado el demandado, nacen a su cargo las cargas de comparecer y contestar, cuyo incumplimiento —a diferencia del deber procesal— no genera una sanción directa pero puede favorecer al actor. De la combinación del cumplimiento o incumplimiento de esas cargas surgen las actitudes omisivas y positivas del demandado que se desarrollan en los capítulos siguientes.
} \\
\bottomrule
\end{tabularx}
\end{table}

% =====================================================
\chapter{La contestación de demanda}
% =====================================================

\section{El plazo para contestar}

El plazo para comparecer y contestar depende del tipo de proceso de que se trate.

\begin{table}[H]
\centering
\caption{Plazos para contestar demanda}
\begin{tabularx}{\textwidth}{l X}
\toprule
\textbf{Proceso} & \textbf{Plazo} \\
\midrule
Ordinario & 15 días hábiles \\
Sumarísimo & 5 días hábiles \\
\bottomrule
\end{tabularx}
\end{table}

\begin{formula}{Art. 156 CPCCN — Plazo de gracia}
El escrito no presentado dentro del horario judicial del día en que venciere el plazo sólo podrá ser entregado válidamente el día hábil inmediato siguiente, durante las dos primeras horas del despacho. Esto significa que quien tiene un plazo de 15 días puede presentar la contestación en las primeras horas hábiles del día 16 (entre las 7:30 y las 9:30 h).
\end{formula}

\textbf{Extensión del plazo por distancia:} si el domicilio del demandado se encuentra a más de 100 km de la sede del tribunal, el plazo se extiende un día cada 200 km o fracción mayor a 100 km. Este plazo es perentorio e improrrogable: si no se contesta en tiempo, se tiene por no presentada la contestación.

\textbf{El plazo es individual:} por más que haya multiplicidad de demandados, cada uno computa el plazo de modo individual, a partir del día siguiente al de la notificación.

\section{Concepto de contestación de demanda}

\begin{definition}{Contestación de demanda}
Es un acto procesal de postulación en el cual el demandado se defiende, pidiéndole a la jurisdicción que rechace la pretensión del actor invocada en la demanda. Es el acto de defensa por excelencia dentro del proceso ordinario, y en él el demandado puede ejercer todas sus defensas de modo amplio.
\end{definition}

\section{Requisitos fiscales y de admisibilidad}

A diferencia de la demanda, al contestar no hay que abonar tasa de justicia (ley 23.898), salvo que el demandado reconvenga. Sin embargo, sí hay que acompañar el bono profesional del Colegio de Abogados en la primera presentación.

\begin{note}
El requisito de mediación previa sólo es exigible al actor; el demandado que advierte que no participó puede pedir al juez que suspenda el proceso y ordene reabrir la instancia de mediación.
\end{note}

\section{Requisitos formales comunes con la demanda}

Tanto la demanda como la contestación deben cumplir con los requisitos formales de todo escrito judicial:

\begin{longtable}{p{0.32\textwidth} p{0.58\textwidth}}
\caption{Requisitos formales de la contestación de demanda} \\
\toprule
\textbf{Requisito} & \textbf{Detalle} \\
\midrule
\endfirsthead
\toprule
\textbf{Requisito} & \textbf{Detalle} \\
\midrule
\endhead
Escrito & Debe hacerse por escrito en idioma nacional. \\
Firma & La firma del letrado y/o de la parte es obligatoria. \\
Representación letrada (Art. 56) & El demandado puede presentarse: (a) por sí, con patrocinio letrado (firman ambos); (b) por apoderado (sólo firma el abogado); o (c) invocando la calidad de gestor (Art. 48) en situaciones urgentes, con plazo de 40 días hábiles para acreditar la representación o que la parte ratifique lo actuado. \\
Denuncia de domicilio real & Debe indicarse el domicilio real del demandado. \\
Constitución de domicilio procesal y electrónico & Es obligatorio constituir domicilio procesal (físico/postal) y electrónico, donde se cursarán todas las notificaciones. \\
\bottomrule
\end{longtable}

\begin{formula}{Art. 48 CPCCN — Gestor procesal}
Cuando no hubiere tiempo para obtener el poder y hubiere urgencia en la realización del acto, se podrá actuar como gestor. El gestor tiene 40 días hábiles para acreditar la representación o que la parte ratifique todo lo actuado. Si no lo hace, se puede decretar la nulidad de lo actuado y se tendrá por no contestada la demanda.
\end{formula}

\section{Presentación digital obligatoria}

No basta con la presentación en formato papel. Dentro de las 24 horas posteriores a la presentación en formato papel, el letrado debe digitalizar esa presentación: escanear el escrito y la documental, y enviarlo electrónicamente al juzgado.

\begin{important}
Si no se cumple con la incorporación digital en el plazo de 24 horas, se cursará una intimación para que se cumpla con la digitalización. Si tampoco se cumple en ese nuevo plazo (24 a 48 horas), se tendrá por no presentado el escrito, con las graves consecuencias que ello implica (se tendría por no contestada la demanda).
\end{important}

Esta exigencia es relativamente nueva y surge de la tendencia hacia el expediente electrónico. El código regula el artículo 120 que, si bien fue pensado para acompañar copias en papel, tiene una aplicación similar: si no se cumple, se puede tener por no presentado el escrito.

\section{Contenido sustancial: artículo 356 CPCCN}

\begin{formula}{Art. 356, inc. 1° CPCCN — Reconocimiento o negación de hechos y documentos}
Al contestar la demanda, el demandado deberá:
\begin{itemize}
    \item Reconocer o negar categóricamente cada uno de los hechos expuestos en la demanda.
    \item Reconocer o negar la autenticidad de los documentos que se le atribuyan y la recepción de las cartas y telegramas a él dirigidos.
\end{itemize}
El silencio, las respuestas evasivas o la negativa meramente general podrán estimarse como reconocimiento de la verdad de los hechos pertinentes y lícitos a que se refieran. En cuanto a los documentos, se los tendrá por reconocidos o recibidos, según el caso.
\end{formula}

Esto implica que el demandado debe:

\begin{table}[H]
\centering
\caption{Consecuencias del silencio y de la negación categórica}
\begin{tabularx}{\textwidth}{p{0.22\textwidth} X X}
\toprule
\textbf{Objeto} & \textbf{Consecuencia del silencio o negativa genérica} & \textbf{Consecuencia de la negación categórica} \\
\midrule
Hechos invocados por el actor & El juez \textbf{puede} tener tácitamente por admitidos los hechos (es una facultad del juez). & El hecho queda controvertido y deberá ser objeto de prueba. \\
Documentos acompañados por el actor & El juez \textbf{debe} tener por reconocidos y auténticos los documentos (es una obligación del juez, no una facultad). & El documento queda desconocido; el actor deberá probar su autenticidad. \\
\bottomrule
\end{tabularx}
\end{table}

\section{Instrumentos públicos vs. instrumentos privados}

Los documentos que acompaña el actor a su demanda pueden ser instrumentos públicos o instrumentos privados. No es la misma eficacia probatoria la de un instrumento privado que la de un instrumento público.

\begin{table}[H]
\centering
\caption{Instrumentos privados vs. instrumentos públicos}
\begin{tabularx}{\textwidth}{p{0.24\textwidth} X X}
\toprule
 & \textbf{Instrumento privado} & \textbf{Instrumento público} \\
\midrule
Ejemplos & Constancia médica, factura, remito de entrega de mercaderías. & Escritura pública otorgada ante escribano. \\
Eficacia probatoria & Menor. No hace plena fe por sí mismo. & Mayor. Hace plena fe por sí mismo, ya que en su confección intervino un oficial público (escribano). \\
Para desconocerlo al contestar (Art. 356) & Basta negar categóricamente su autenticidad, uno a uno. Ej.: ``Niega la autenticidad en forma y contenido de la constancia médica emitida por el Sanatorio Los Arcos\ldots'' & Se lo debe impugnar al contestar, pero ello por sí solo no basta para restarle eficacia probatoria. \\
Para quitarle eficacia probatoria definitivamente & Con la negación categórica en la contestación es suficiente. El actor deberá probar su autenticidad. & Debe promoverse un incidente de redargución de falsedad (Art. 395 CPCCN) dentro de los 15 días de haberlo impugnado. Si no se promueve, el documento sigue teniendo plena eficacia probatoria. \\
\bottomrule
\end{tabularx}
\end{table}

\begin{formula}{Art. 395 CPCCN — Incidente de redargución de falsedad}
La redargución de falsedad de un instrumento público tramitará por incidente que deberá promoverse dentro del plazo de 15 días desde que se tachó de falso el documento. El incidente no suspenderá el proceso principal salvo que el juez lo disponga.
\end{formula}

\begin{example}{Reconocimiento de firma en instrumento privado}
Si el actor acompaña un documento cuya firma se atribuye al demandado, y el demandado reconoce dicha firma, se tendrá por auténtico todo el contenido del documento.
\end{example}

\section{Art. 356 inc. 2°: relación de los hechos del demandado}

\begin{formula}{Art. 356, inc. 2° CPCCN — Hechos que alega el demandado}
El demandado deberá especificar con claridad los hechos que alegare como fundamento de su defensa.
\end{formula}

El demandado debe dar su versión de los hechos, así como el actor lo hizo en la demanda según el artículo 330. Describirá los hechos en virtud de los cuales solicita a la jurisdicción que rechace la pretensión del actor.

\begin{example}{Versión del demandado en diferentes procesos}
\begin{enumerate}
    \item Mala praxis médica: el médico demandado dará su versión de los hechos para tratar de demostrar que no existió el acto o la conducta tachada de mala praxis.
    \item Accidente de tránsito: el demandado relatará cuál es la forma en la que realmente ocurrió el accidente y por qué no es responsable del mismo.
\end{enumerate}
\end{example}

\section{Ofrecimiento de prueba al contestar (Art. 333 CPCCN)}

\begin{formula}{Art. 333 CPCCN — Documentos y ofrecimiento de prueba}
Con la contestación de demanda, la parte demandada deberá acompañar la prueba documental que estuviese en su poder y ofrecer todos los demás medios de prueba de que intentare valerse.
\end{formula}

La oportunidad para ofrecer prueba es, para ambas partes, en los escritos introductorios: en la demanda para el actor, y en la contestación de demanda para el demandado, quien ofrecerá todas las pruebas de las que intente valerse.

\begin{table}[H]
\centering
\caption{Requisitos según el tipo de prueba ofrecida}
\begin{tabularx}{\textwidth}{p{0.22\textwidth} X}
\toprule
\textbf{Tipo de prueba} & \textbf{Requisitos al ofrecerla} \\
\midrule
Documental & Si está en poder del demandado, debe acompañarla al escrito de contestación. \\
Testimonial & Debe individualizar a los testigos con nombre, apellido, documento, profesión y domicilio. Debe indicar a qué fines ofrece cada testigo y qué extremos intenta acreditar con su declaración. \\
Pericial & Debe indicar el tipo de perito (ej.: psicológica, médica, contable). También debe indicar los puntos de pericia: el cuestionario que pretende que el perito designado en la etapa probatoria responda en su dictamen. \\
\bottomrule
\end{tabularx}
\end{table}

\begin{example}{Ofrecimiento de testigo en accidente de tránsito}
Se ofrece como testigo a Juan González, DNI\ldots, domiciliado en\ldots, en carácter de testigo presencial del accidente, a fin de que declare respecto de la forma de ocurrencia del mismo. En un caso de mala praxis de traumatología, podría ofrecerse como testigo técnico a un reconocido traumatólogo para que brinde un testimonio técnico sobre la conducta médica en cuestión.
\end{example}

\section{El derecho invocado y la petición (remisión al Art. 330)}

El artículo 356 señala que también deben cumplirse las disposiciones del artículo 330 que resulten aplicables:

\begin{table}[H]
\centering
\caption{Requisitos del Art. 330 aplicables a la contestación}
\begin{tabularx}{\textwidth}{p{0.30\textwidth} X}
\toprule
\textbf{Requisito del Art. 330} & \textbf{Aplicación a la contestación} \\
\midrule
Invocación del derecho & El demandado debe invocar de modo concreto y claro el derecho en el cual funda su defensa. \\
Petición clara y concreta & La petición del demandado será: que se lo tenga por presentado, por parte, por contestada la demanda, por ofrecida la prueba, y que oportunamente se rechace la demanda y la pretensión del actor. \\
\bottomrule
\end{tabularx}
\end{table}

\begin{important}
\textbf{Principio iura novit curia:} el juez conoce el derecho. Aunque las partes deben invocar el derecho en sus escritos, el juez puede apartarse de esa invocación y aplicar el derecho que considera correcto al momento de dictar sentencia. Esto es absolutamente válido.
\end{important}

\section{Efectos de la contestación de demanda}

\begin{longtable}{p{0.32\textwidth} p{0.58\textwidth}}
\caption{Efectos de la contestación de demanda} \\
\toprule
\textbf{Efecto} & \textbf{Descripción} \\
\midrule
\endfirsthead
\toprule
\textbf{Efecto} & \textbf{Descripción} \\
\midrule
\endhead
Cierre de la etapa introductoria & Con la contestación se clausura la etapa introductoria del proceso. \\
Fijación de la competencia & Si el demandado no cuestionó la competencia, queda establecido que ese juez entenderá a lo largo de todo el proceso. \\
Imposibilidad de recusar sin causa & La oportunidad para recusar sin causa al juez era en la contestación (Art. 14 CPCCN). Si no se ejerció, ya no puede hacerse. \\
Preclusión de excepciones & Ya no podrán oponerse excepciones que no fueron planteadas en la contestación. \\
Preclusión de reconvención & Si no se reconvino al contestar, no puede hacerse después. \\
Fijación de la cuestión litigiosa & Con la demanda y la contestación queda fijada la cuestión litigiosa: se identifican los hechos afirmados por cada parte, los hechos controvertidos y los hechos conducentes. Esto determina cuáles serán objeto de prueba. \\
\bottomrule
\end{longtable}

\section{Traslado de la documentación adjuntada por el demandado}

Si el demandado acompañó prueba documental, en la resolución que dicte el juez al tener por presentada la contestación de demanda también se resolverá dar traslado de esa documentación al actor para que, en el plazo de cinco días, se pronuncie respecto de ella en los términos del artículo 356, inciso 1°.

Así, se mantienen el principio de bilateralidad y la igualdad de las partes ante la jurisdicción: así como el demandado tuvo que pronunciarse sobre la documentación del actor, el actor debe pronunciarse sobre la documentación del demandado. Este traslado se efectúa por cédula.

\section*{Cuadro Cornell del capítulo}
\addcontentsline{toc}{section}{Cuadro Cornell del capítulo}

\begin{table}[H]
\centering
\begin{tabularx}{\textwidth}{
    >{\raggedright\arraybackslash}p{0.30\textwidth}
    >{\raggedright\arraybackslash}X
}
\toprule
\textbf{Preguntas / palabras clave} & \textbf{Notas / respuestas} \\
\midrule
\textbf{¿Cuál es el plazo para contestar en el proceso ordinario y en el sumarísimo?} & 15 días hábiles en el ordinario y 5 días hábiles en el sumarísimo, contados desde el día siguiente a la notificación, con el plazo de gracia del Art. 156 CPCCN. \\
\textbf{¿Qué debe hacer el demandado con los hechos y los documentos invocados por el actor (Art. 356 inc. 1°)?} & Debe reconocer o negar categóricamente cada hecho y cada documento. El silencio o la negativa genérica facultan al juez a tener por admitidos los hechos, y lo obligan a tener por reconocidos los documentos. \\
\textbf{¿Cómo se impugna un instrumento público para quitarle eficacia probatoria?} & Debe impugnarse al contestar y, además, promoverse un incidente de redargución de falsedad (Art. 395 CPCCN) dentro de los 15 días de la impugnación; de lo contrario, conserva plena eficacia probatoria. \\
\textbf{¿Cuáles son los principales efectos de la contestación de demanda?} & Cierra la etapa introductoria, fija la competencia, impide la recusación sin causa, produce la preclusión de excepciones y de la reconvención, y fija la cuestión litigiosa. \\
\midrule
\multicolumn{2}{p{\textwidth}}{
\textbf{Resumen:} La contestación de demanda es el acto de defensa por excelencia del demandado, regulado por el Art. 356 CPCCN, que le impone reconocer o negar categóricamente hechos y documentos, con consecuencias distintas según guarde silencio o niegue expresamente. Junto con la contestación debe ofrecerse toda la prueba de la que se intente valer y acompañarse la documental disponible. La contestación produce efectos preclusivos relevantes: cierra la etapa introductoria y fija definitivamente la competencia, la posibilidad de recusar sin causa, las excepciones y la reconvención.
} \\
\bottomrule
\end{tabularx}
\end{table}

% =====================================================
\chapter{Documentos nuevos y hechos no invocados}
% =====================================================

\section{La regla general: oportunidad para ofrecer prueba}

En principio, la única oportunidad para ofrecer prueba documental es la de los escritos introductorios. Sin embargo, el código regula algunas situaciones puntuales en las cuales se admite incorporar más prueba documental o alegar hechos adicionales.

\section{Documentos nuevos (Art. 335 CPCCN)}

\begin{formula}{Art. 335 CPCCN — Documentos posteriores o desconocidos}
Después de interpuesta la demanda o contestada ésta, no se admitirán al actor ni al demandado, respectivamente, otros documentos que los de fecha posterior a aquéllos, o los que, siendo de fecha anterior, se juren no haber tenido antes conocimiento de ellos.
\end{formula}

Se pueden incorporar documentos con posterioridad a la demanda o a la contestación en dos supuestos:

\begin{table}[H]
\centering
\caption{Supuestos de incorporación de documentos nuevos}
\begin{tabularx}{\textwidth}{p{0.32\textwidth} X}
\toprule
\textbf{Supuesto} & \textbf{Descripción} \\
\midrule
Documentos de fecha posterior & Son documentos emitidos con fecha posterior a la promoción de la demanda o a la contestación. Es lógico que no hubieran podido acompañarse antes. \\
Documentos de fecha anterior desconocidos & Documentos con fecha anterior al inicio de la demanda o la contestación, pero cuya existencia las partes desconocían. Deben declarar bajo juramento que no tenían conocimiento de su existencia al momento de promover la demanda o contestar. \\
\bottomrule
\end{tabularx}
\end{table}

Presentado el documento, el principio de bilateralidad impone correr traslado a la otra parte por cinco días para que se pronuncie sobre su reconocimiento o desconocimiento en los términos del artículo 356, inciso 1°.

En primera instancia, los documentos pueden incorporarse hasta que los autos se encuentren para sentencia (es decir, hasta que el juez dicte el decreto ``autos para sentencia''). Una vez dictado ese decreto, ya no se puede realizar actividad procesal alguna, y el expediente espera el dictado de la sentencia. Si los documentos surgieran después, la oportunidad para incorporarlos es ante la Cámara, al fundar los recursos de apelación (Art. 260 CPCCN).

\begin{example}{Documentos de fecha posterior (obra lindera)}
Se promueve una demanda por los daños generados a un inmueble por una obra en construcción lindera. En ese momento, los daños aún no están reparados y el actor reclama en base a un presupuesto de reparaciones. A los cuatro meses del inicio del proceso, las reparaciones se realizan y se pagan, por lo que el actor cuenta ahora con una factura (documento de fecha posterior) que acredita lo que realmente pagó. Esta factura puede incorporarse al proceso.
\end{example}

\section{Hechos no invocados (Art. 334 CPCCN)}

\begin{formula}{Art. 334 CPCCN — Hechos no invocados en los escritos de constitución del proceso}
Si en el responde de la demanda se alegaren hechos no invocados en ésta, el actor podrá, en el plazo de cinco días contado desde la notificación del traslado mencionado en el artículo 150, ampliar o modificar la prueba ofrecida.
\end{formula}

El demandado puede introducir hechos que no fueron invocados por el actor en su relato, pero que claramente se relacionan con él y están vinculados al desarrollo de la relación jurídica entre el actor y el demandado. Por alguna razón —por omisión o por una cuestión de estrategia— el actor no los mencionó, y es el demandado quien los introduce al contestar la demanda.

Si el demandado introduce hechos no invocados en la contestación, el actor tiene cinco días desde que se le notifica la resolución que tuvo por contestada la demanda para agregar documentos y ofrecer pruebas vinculadas a esos hechos no invocados.

\section{Hechos nuevos (Art. 365 CPCCN)}

\begin{formula}{Art. 365 CPCCN — Hechos nuevos}
Cuando con posterioridad a la contestación de la demanda o reconvención ocurriese o llegase a conocimiento de las partes algún hecho que tuviese relación con la cuestión que se ventila, podrán alegarlo hasta cinco días después de notificada la audiencia prevista en el artículo 360.
\end{formula}

\begin{example}{Hecho nuevo en juicio de mala praxis}
Se promueve una demanda por mala praxis médica que generó un daño cardíaco al actor. Luego de iniciada la demanda, el actor debe ser intervenido quirúrgicamente nuevamente por la misma causa. Esta nueva intervención quirúrgica es un hecho nuevo sobreviniente, que puede ser invocado como tal en el proceso.
\end{example}

A continuación se presenta un cuadro comparativo entre las tres figuras estudiadas en este capítulo:

\begin{table}[H]
\centering
\caption{Documentos nuevos, hechos no invocados y hechos nuevos}
\begin{tabularx}{\textwidth}{p{0.26\textwidth} X X X}
\toprule
 & \textbf{Documentos nuevos (Art. 335)} & \textbf{Hechos no invocados (Art. 334)} & \textbf{Hechos nuevos (Art. 365)} \\
\midrule
Oportunidad en 1\textsuperscript{a} instancia & Hasta los autos para sentencia. & Al contestar demanda (el demandado); el actor responde dentro de los 5 días. & Hasta 5 días después de notificada la audiencia del Art. 360. \\
Oportunidad en Cámara & Sí (Art. 260 CPCCN). & No corresponde. & No corresponde. \\
Característica principal & El documento no existía o era desconocido al momento de la demanda/contestación. & El hecho existía pero no fue mencionado por el actor; el demandado lo introduce. & El hecho ocurrió o llegó a conocimiento de las partes con posterioridad a la contestación. \\
\bottomrule
\end{tabularx}
\end{table}

\section*{Cuadro Cornell del capítulo}
\addcontentsline{toc}{section}{Cuadro Cornell del capítulo}

\begin{table}[H]
\centering
\begin{tabularx}{\textwidth}{
    >{\raggedright\arraybackslash}p{0.30\textwidth}
    >{\raggedright\arraybackslash}X
}
\toprule
\textbf{Preguntas / palabras clave} & \textbf{Notas / respuestas} \\
\midrule
\textbf{¿En qué dos supuestos pueden incorporarse documentos nuevos (Art. 335)?} & Cuando son de fecha posterior a la demanda/contestación, o cuando siendo de fecha anterior, la parte jura no haber tenido conocimiento de ellos antes. \\
\textbf{¿Hasta cuándo pueden incorporarse documentos en primera instancia?} & Hasta que se dicte el decreto de ``autos para sentencia''; luego sólo pueden incorporarse ante la Cámara, al fundar los recursos de apelación (Art. 260 CPCCN). \\
\textbf{¿Qué son los hechos no invocados (Art. 334) y quién los introduce?} & Son hechos vinculados a la relación jurídica que el actor no mencionó en la demanda; el demandado los introduce al contestar, y el actor cuenta con 5 días para ofrecer prueba sobre ellos. \\
\textbf{¿Qué diferencia a los hechos nuevos (Art. 365) de los hechos no invocados?} & Los hechos nuevos ocurren o se conocen después de contestada la demanda, y pueden alegarse hasta 5 días después de notificada la audiencia del Art. 360; los hechos no invocados ya existían al momento de demandar. \\
\midrule
\multicolumn{2}{p{\textwidth}}{
\textbf{Resumen:} Aunque la regla general es que la prueba se ofrece únicamente en los escritos introductorios, el CPCCN prevé tres excepciones puntuales: los documentos nuevos (Art. 335), que pueden incorporarse hasta los autos para sentencia en primera instancia o al fundar la apelación en Cámara; los hechos no invocados (Art. 334), que el demandado introduce al contestar y que el actor puede rebatir dentro de los cinco días; y los hechos nuevos (Art. 365), sobrevinientes a la contestación, que pueden alegarse hasta cinco días después de notificada la audiencia del Art. 360.
} \\
\bottomrule
\end{tabularx}
\end{table}

% =====================================================
\chapter{El allanamiento}
% =====================================================

\section{Concepto y naturaleza jurídica}

\begin{definition}{Allanamiento}
Es el acto procesal en virtud del cual, de modo unilateral (sin necesitar la conformidad de la parte actora), el demandado se somete a la pretensión que el actor invocó en su demanda. No implica reconocer como verdaderos los hechos y el derecho invocados por el actor; simplemente el demandado se somete a la pretensión.
\end{definition}

El allanamiento es uno de los modos anormales (o excepcionales) de terminar el proceso, regulado en el artículo 307 del CPCCN. Se denomina ``anormal'' porque el modo habitual de terminar el proceso es con el dictado de la sentencia; el allanamiento lo hace terminar de una manera diferente.

\section{Tipos de allanamiento}

\begin{table}[H]
\centering
\caption{Tipos de allanamiento}
\begin{tabularx}{\textwidth}{p{0.18\textwidth} X}
\toprule
\textbf{Tipo} & \textbf{Descripción} \\
\midrule
Total & El demandado se somete completamente a la pretensión del actor. \\
Parcial & El demandado se somete sólo a una parte de la pretensión del actor. \\
Expreso & El demandado manifiesta de modo explícito en su escrito que se allana y se somete a la pretensión. \\
Tácito & Se infiere del comportamiento del demandado. Ej.: en un proceso de desalojo, el demandado notificado deposita las llaves del inmueble en el expediente dentro del plazo para contestar. \\
\bottomrule
\end{tabularx}
\end{table}

\begin{example}{Allanamiento tácito en proceso de desalojo}
En un proceso de desalojo, el demandado —al ser notificado del inicio de la demanda y dentro del plazo para contestar— se presenta y deposita o consigna las llaves del inmueble en el expediente. Esta conducta tácita implica un allanamiento: el demandado se somete a la pretensión del actor y devuelve la llave, desocupando el inmueble objeto del proceso.
\end{example}

\section{Oportunidad del allanamiento}

El demandado puede allanarse en cualquier instancia del proceso, en tanto todavía no se haya dictado sentencia. Una de las oportunidades en que puede hacerlo —vinculada con lo estudiado en esta unidad— es justamente la de contestar la demanda.

Las razones que pueden llevar al demandado a allanarse son variadas:

\begin{table}[H]
\centering
\caption{Razones para allanarse}
\begin{tabularx}{\textwidth}{p{0.30\textwidth} X}
\toprule
\textbf{Razón} & \textbf{Descripción} \\
\midrule
Conciencia del resultado probable & El demandado entiende que la sentencia probablemente hará lugar a la demanda. \\
Evitar mayores costas & Allanarse antes evita las costas que generaría tramitar todas las etapas del proceso. \\
Evitar la incertidumbre & No queda ``atado'' al resultado de un proceso que puede durar años en el caso ordinario. \\
\bottomrule
\end{tabularx}
\end{table}

\section*{Cuadro Cornell del capítulo}
\addcontentsline{toc}{section}{Cuadro Cornell del capítulo}

\begin{table}[H]
\centering
\begin{tabularx}{\textwidth}{
    >{\raggedright\arraybackslash}p{0.30\textwidth}
    >{\raggedright\arraybackslash}X
}
\toprule
\textbf{Preguntas / palabras clave} & \textbf{Notas / respuestas} \\
\midrule
\textbf{¿Qué es el allanamiento y qué implica respecto de los hechos y el derecho invocados por el actor?} & Es el acto unilateral por el cual el demandado se somete a la pretensión del actor, sin que ello implique reconocer como verdaderos los hechos y el derecho invocados. \\
\textbf{¿Por qué se considera un modo anormal de terminar el proceso?} & Porque el modo habitual de terminación es el dictado de la sentencia, y el allanamiento pone fin al proceso de una manera diferente (Art. 307 CPCCN). \\
\textbf{¿Qué tipos de allanamiento existen?} & Total o parcial, según el alcance de la sumisión; y expreso o tácito, según cómo se manifieste. \\
\textbf{¿Hasta qué momento puede allanarse el demandado?} & En cualquier instancia del proceso, siempre que todavía no se haya dictado sentencia. \\
\midrule
\multicolumn{2}{p{\textwidth}}{
\textbf{Resumen:} El allanamiento es un modo anormal y unilateral de terminación del proceso (Art. 307 CPCCN) por el cual el demandado se somete a la pretensión del actor, sin reconocer necesariamente los hechos ni el derecho invocados. Puede ser total o parcial y expreso o tácito, y puede ejercerse en cualquier momento anterior a la sentencia, siendo la contestación de demanda una de sus oportunidades típicas. Entre sus razones más comunes están la conciencia del resultado probable del pleito, el ahorro de costas y la eliminación de la incertidumbre del proceso.
} \\
\bottomrule
\end{tabularx}
\end{table}

% =====================================================
\chapter{La reconvención}
% =====================================================

\section{Concepto}

\begin{definition}{Reconvención}
Es la introducción, dentro del mismo proceso ordinario, de una demanda que el demandado dirige contra el actor. Contiene una pretensión autónoma y nueva que, en lugar de venir del actor (como en la demanda original), proviene del demandado hacia el actor, para que sea éste quien la satisfaga. El fundamento de su regulación es la economía procesal: evitar el inicio de otro proceso separado para debatir una cuestión vinculada.
\end{definition}

\section{Doble calidad de las partes}

El actor pasa a ser reconvenido y el demandado, a su vez, reconviniente. Ambos sujetos tienen, entonces, la calidad de actores y de demandados dentro del mismo proceso.

El juez, al dictar sentencia, analizará tanto la pretensión del actor en la demanda como la pretensión del demandado en la reconvención, dictando una única sentencia para ambas.

\section{Oportunidad y plazo}

La oportunidad para reconvenir es la misma que la del plazo para contestar demanda: en el proceso ordinario, dentro de los 15 días para contestar demanda.

\begin{formula}{Art. 357 y ss. CPCCN — Reconvención}
La reconvención se deducirá en el mismo escrito de contestación, observándose los requisitos aplicables a la demanda. De la reconvención se dará traslado al actor por el plazo de 15 días.
\end{formula}

\section{Trámite procesal de la reconvención}

\begin{longtable}{p{0.32\textwidth} p{0.58\textwidth}}
\caption{Trámite procesal de la reconvención} \\
\toprule
\textbf{Paso} & \textbf{Descripción} \\
\midrule
\endfirsthead
\toprule
\textbf{Paso} & \textbf{Descripción} \\
\midrule
\endhead
1. Demanda & El actor promueve demanda. Se notifica al demandado y corre el plazo de 15 días. \\
2. Contestación + reconvención & El demandado contesta la demanda y, en el mismo escrito, reconviene. El acto es simultáneo. \\
3. Traslado de la reconvención al actor & El juez da traslado al actor (ahora reconvenido) por 15 días para que conteste la reconvención y oponga defensas (incluyendo las excepciones de los Arts. 346, 347 y 348). \\
4. Cierre de la etapa introductoria & Concluidos los traslados, la etapa introductoria se cierra con cuatro escritos: demanda, contestación, reconvención y contestación de la reconvención. \\
5. Audiencia preliminar (Art. 360) & El juez fija y celebra la audiencia, teniendo en cuenta todos los escritos introductorios para identificar hechos controvertidos, conducentes y prueba a producir. \\
6. Sentencia única & El juez dicta una única sentencia que resuelve tanto la pretensión del actor como la del demandado reconviniente. \\
\bottomrule
\end{longtable}

\section{Requisitos de la reconvención}

\begin{longtable}{p{0.32\textwidth} p{0.58\textwidth}}
\caption{Requisitos de la reconvención} \\
\toprule
\textbf{Requisito} & \textbf{Descripción} \\
\midrule
\endfirsthead
\toprule
\textbf{Requisito} & \textbf{Descripción} \\
\midrule
\endhead
Mismos requisitos que la demanda (Art. 330) & Al ser una demanda, debe cumplir con los requisitos formales y sustanciales de la demanda. \\
Requisitos fiscales & El demandado debe abonar la tasa de justicia vinculada al monto reclamado en la reconvención (Ley 23.898). Si no cuenta con recursos, puede promover un beneficio de litigar sin gastos. \\
Identidad de sujetos & Debe tratarse de los mismos sujetos (actor y demandado en el proceso principal). \\
Mismo tipo de proceso & Ambas pretensiones deben poder tramitar en el mismo tipo de proceso (ej.: ambas en proceso ordinario). \\
Proceso en trámite con traslado notificado & Debe haber un proceso en trámite y el traslado de la demanda debe haber sido notificado. No puede reconvenirse en diligencias preliminares ni en prueba anticipada. \\
Relación entre pretensiones & Las pretensiones deben estar relacionadas: deben derivar de la misma relación jurídica, o ser conexas entre sí. \\
\bottomrule
\end{longtable}

\section{Alternativa a la reconvención: acumulación de procesos}

Si el demandado no reconviene al contestar, puede iniciar luego un nuevo proceso, siempre que no haya operado la prescripción. Pero entonces habrá dos procesos vinculados, en los cuales el actor de uno es demandado en el otro. Para evitar sentencias contradictorias, estos procesos deben acumularse.

\begin{important}
\textbf{Acumulación de procesos:} ambos expedientes tramitarán ante el mismo juez, por aplicación del principio de prevención: el juez que entiende en el expediente en el que primero se notificó la demanda. El otro expediente se remite a ese juzgado. Cuando ambos estén en condiciones de dictar sentencia, el juez dictará una única sentencia para ambos procesos acumulados.
\end{important}

\section{Ejemplos de reconvención}

\begin{example}{Contrato de compraventa}
Las partes celebraron un contrato de compraventa. El actor (comprador) promueve demanda reclamando la entrega de la cosa adquirida. El demandado (vendedor), además de contestar, puede reconvenir reclamando al actor el pago del saldo de precio. Ambas pretensiones derivan de la misma relación jurídica: el contrato de compraventa.
\end{example}

\begin{example}{Accidente de tránsito}
El actor demanda al demandado reclamando una indemnización por daños y perjuicios, argumentando que el demandado cruzó en rojo y causó el accidente. El demandado contesta, dando su propia versión: fue el actor quien cruzó en rojo. Además de contestar defensivamente, el demandado reconviene, reclamando a su vez una indemnización por los daños que sufrió en el mismo accidente. Las pretensiones son conexas: ambas derivan del mismo accidente de tránsito.
\end{example}

\begin{example}{Cuota alimentaria (fuero de familia)}
Los padres de tres menores acordaron, luego del divorcio, que el padre abonaría \$10.000 mensuales en concepto de cuota alimentaria. Cinco años después, la madre promueve un expediente solicitando el aumento de esa cuota. El demandado (padre), al contestar, puede reconvenir solicitando la reducción de la cuota, argumentando que las circunstancias cambiaron: la madre, que antes estaba desempleada, ahora trabaja como gerente en una empresa multinacional.
\end{example}

\section*{Cuadro Cornell del capítulo}
\addcontentsline{toc}{section}{Cuadro Cornell del capítulo}

\begin{table}[H]
\centering
\begin{tabularx}{\textwidth}{
    >{\raggedright\arraybackslash}p{0.30\textwidth}
    >{\raggedright\arraybackslash}X
}
\toprule
\textbf{Preguntas / palabras clave} & \textbf{Notas / respuestas} \\
\midrule
\textbf{¿Qué es la reconvención y cuál es su fundamento?} & Es la demanda que el demandado dirige contra el actor dentro del mismo proceso, ejerciendo una pretensión autónoma. Su fundamento es la economía procesal: evitar dos procesos paralelos y el riesgo de sentencias contradictorias. \\
\textbf{¿Cuál es la oportunidad y el plazo para reconvenir en el proceso ordinario?} & Debe deducirse en el mismo escrito de contestación, dentro de los 15 días previstos para contestar demanda. \\
\textbf{¿Qué requisito deben cumplir las pretensiones del actor y del reconviniente entre sí?} & Deben estar relacionadas: derivar de la misma relación jurídica o ser conexas entre sí. \\
\textbf{¿Qué ocurre si el demandado no reconviene al contestar?} & Puede iniciar luego un proceso separado, si no operó la prescripción, pero ambos expedientes deberán acumularse ante el mismo juez para evitar sentencias contradictorias. \\
\midrule
\multicolumn{2}{p{\textwidth}}{
\textbf{Resumen:} La reconvención es la demanda que el demandado dirige contra el actor dentro del mismo proceso, fundada en la economía procesal. Debe deducirse junto con la contestación, dentro del mismo plazo, y requiere identidad de sujetos, mismo tipo de proceso y relación o conexidad entre las pretensiones. De su traslado y contestación surgen cuatro escritos introductorios (demanda, contestación, reconvención y contestación de la reconvención), sobre los cuales el juez dictará una única sentencia. Si el demandado no reconviene, puede iniciar un proceso separado, pero ambos expedientes deberán acumularse para evitar decisiones contradictorias.
} \\
\bottomrule
\end{tabularx}
\end{table}

% =====================================================
\chapter{La rebeldía}
% =====================================================

\section{Concepto y diferencia con la simple incomparecencia}

\begin{definition}{Rebeldía}
Situación procesal que se configura cuando el demandado, debidamente notificado de la demanda, no comparece y por ende tampoco contesta (Arts. 59 y ss. CPCCN). También puede darse cuando una parte que estaba presentada en el proceso lo abandona en el transcurso del mismo.
\end{definition}

\begin{table}[H]
\centering
\caption{Comparece y no contesta vs. rebeldía}
\begin{tabularx}{\textwidth}{p{0.30\textwidth} X X}
\toprule
 & \textbf{Comparece y no contesta} & \textbf{No comparece y no contesta (rebeldía)} \\
\midrule
Cumple la carga de comparecer & Sí & No \\
Cumple la carga de contestar & No & No \\
Notificaciones & Por cédula al domicilio electrónico constituido. & Por ministerio de ley (martes y viernes en sede del tribunal), hasta que se decrete la rebeldía y se la notifique por cédula. \\
Puede ejercer actos procesales futuros & Sí, todos los pertinentes desde ese momento. & No puede ejercer actos de etapas anteriores a su presentación. \\
\bottomrule
\end{tabularx}
\end{table}

\section{Declaración de la rebeldía}

La rebeldía se decreta a pedido del actor, no de oficio. Si el actor no solicita la rebeldía del demandado, el proceso continúa y todas las resoluciones que se vayan dictando se entienden notificadas a ese demandado por ministerio de ley.

El actor presenta un escrito solicitando la declaración de rebeldía, invocando los artículos 59 y siguientes del CPCCN, y señalando que el plazo de 15 días venció sin que el demandado se haya presentado.

El juez dicta una resolución (providencia) declarando al demandado rebelde y ordenando que las resoluciones futuras se tendrán por notificadas por ministerio de ley. Esta resolución se notifica por cédula al domicilio real del demandado, al igual que la demanda.

Una vez dictada la sentencia de primera instancia, ella también debe notificarse al rebelde por cédula al domicilio real.

\section{Efectos de la rebeldía}

\begin{longtable}{p{0.32\textwidth} p{0.58\textwidth}}
\caption{Efectos de la rebeldía} \\
\toprule
\textbf{Efecto} & \textbf{Descripción} \\
\midrule
\endfirsthead
\toprule
\textbf{Efecto} & \textbf{Descripción} \\
\midrule
\endhead
Proceso continúa normalmente & La rebeldía no implica que los hechos del actor se tengan automáticamente por ciertos ni que se dicte sentencia sin más. El proceso continúa con todas sus etapas. \\
Silencio sobre hechos y documentos (Art. 356 inc. 1°) & Al no contestar, el rebelde guardó silencio respecto de los hechos y documentos del actor. El juez puede tener por admitidos los hechos y debe tener por reconocidos los documentos. \\
Presunción de verdad & El estado de rebeldía implica una presunción de verdad de los hechos invocados por el actor, que el juez debe considerar en caso de duda. \\
Medidas cautelares (Art. 63 y Art. 212) & El actor puede solicitar medidas cautelares para asegurar el efectivo cumplimiento de la sentencia. La rebeldía permite presumir la verosimilitud del derecho y el peligro en la demora. El Art. 212 también permite pedir embargo preventivo. \\
Notificaciones por ministerio de ley & Desde la declaración de rebeldía hasta la sentencia, las resoluciones se notifican al rebelde por ministerio de ley (martes y viernes). \\
\bottomrule
\end{longtable}

\begin{formula}{Art. 63 CPCCN — Medidas cautelares en la rebeldía}
El rebelde podrá pedir la revocación de las medidas cautelares decretadas, si acreditare que la rebeldía fue motivada en impedimento justificado que no pudo hacer saber al tribunal. En cuanto al fondo del asunto, el juez apreciará en la sentencia los efectos de la rebeldía como elemento de convicción corroborante de las pruebas producidas.
\end{formula}

\section{Cesación de la rebeldía: el rebelde que se presenta}

Si el rebelde se presenta en alguna instancia posterior del proceso, puede pedir que cese su rebeldía. Los actos que no pudo realizar en las etapas anteriores ya no podrá realizarlos —no podrá contestar demanda, ni oponer excepciones, ni ofrecer prueba de esa etapa— pero sí podrá, a partir de ese momento, realizar todos los actos procesales pertinentes en las distintas etapas siguientes.

\begin{table}[H]
\centering
\caption{Consecuencias de la presentación del rebelde según el momento procesal}
\begin{tabularx}{\textwidth}{p{0.30\textwidth} X}
\toprule
\textbf{Momento} & \textbf{Consecuencia} \\
\midrule
Antes de la etapa probatoria & No puede contestar demanda ni ofrecer prueba; puede controlar la prueba del actor y ejercer todos los actos futuros. \\
Durante la etapa probatoria & Puede impugnar declaraciones testimoniales, dictámenes periciales, y ejercer los demás actos procesales pertinentes. \\
Medidas cautelares & En principio se mantienen. El rebelde puede pedir su modificación, sustitución o disminución. Si demuestra que su rebeldía fue forzada por causas insuperables (ej.: estaba internado en coma), el juez puede levantarlas. \\
\bottomrule
\end{tabularx}
\end{table}

\begin{example}{Levantamiento de medidas cautelares por causa justificada}
El demandado alega, al presentarse, que se encontraba internado en estado de coma en un sanatorio durante todo el plazo para contestar la demanda. El juez valorará si esa circunstancia justifica la rebeldía y, de considerarlo así, podrá levantar las medidas cautelares trabadas.
\end{example}

\section{Excepción de prescripción para el rebelde (Art. 346 CPCCN)}

\begin{formula}{Art. 346 CPCCN — Excepciones y rebelde}
El rebelde sólo podrá oponer la excepción de prescripción con posterioridad al plazo para oponer excepciones, siempre que justifique haber incurrido en rebeldía por causas que no hayan estado a su alcance superar.
\end{formula}

El rebelde puede, con posterioridad, oponer la excepción de prescripción, siempre que, según la valoración que el juez haga del caso, demuestre que cayó en ese estado de rebeldía por cuestiones insalvables.

\section{Nulidad de la notificación y retroacción del proceso}

Si el demandado no fue notificado debidamente —por ejemplo, si la cédula fue dirigida a un domicilio incorrecto, que no era su residencia habitual ni el registrado en los registros públicos—, puede presentarse y plantear la nulidad de esa notificación. Si el juez hace lugar a la nulidad, se producen las siguientes consecuencias:

\begin{table}[H]
\centering
\caption{Consecuencias de la nulidad de la notificación}
\begin{tabularx}{\textwidth}{p{0.30\textwidth} X}
\toprule
\textbf{Consecuencia} & \textbf{Descripción} \\
\midrule
Retroacción del proceso & Se considera que todos los actos realizados con posterioridad al acto nulo (la notificación) también son nulos; el proceso se retrotrae. \\
Nueva notificación & Debe volver a notificarse debidamente al demandado del traslado de la demanda. \\
El demandado puede contestar & Debidamente notificado, el demandado puede ejercer todos sus derechos, incluyendo contestar la demanda. \\
Levantamiento de medidas cautelares & Si se habían trabado medidas cautelares y prospera la nulidad, también se levantan. \\
\bottomrule
\end{tabularx}
\end{table}

\section*{Cuadro Cornell del capítulo}
\addcontentsline{toc}{section}{Cuadro Cornell del capítulo}

\begin{table}[H]
\centering
\begin{tabularx}{\textwidth}{
    >{\raggedright\arraybackslash}p{0.30\textwidth}
    >{\raggedright\arraybackslash}X
}
\toprule
\textbf{Preguntas / palabras clave} & \textbf{Notas / respuestas} \\
\midrule
\textbf{¿Qué es la rebeldía y cómo se declara?} & Es la situación en que incurre el demandado debidamente notificado que no comparece y no contesta. Se declara a pedido del actor (no de oficio), una vez vencido el plazo para contestar, y la resolución se notifica por cédula al domicilio real. \\
\textbf{¿Cuáles son los principales efectos de la rebeldía?} & El proceso continúa normalmente; se aplica el Art. 356 inc. 1° respecto de hechos y documentos; genera presunción de verdad de los hechos del actor; habilita medidas cautelares; y las resoluciones se notifican por ministerio de ley. \\
\textbf{¿Qué puede hacer el rebelde que se presenta tardíamente?} & Puede ejercer todos los actos procesales desde su presentación en adelante, pero no los de etapas ya precluidas (contestar, ofrecer prueba, oponer excepciones). Puede pedir la modificación o el levantamiento de medidas cautelares si justifica causa insuperable. \\
\textbf{¿Qué ocurre si se declara la nulidad de la notificación de la demanda?} & El proceso se retrotrae: se anulan los actos posteriores, debe notificarse nuevamente al demandado, éste puede contestar y se levantan las medidas cautelares que se hubieran trabado. \\
\midrule
\multicolumn{2}{p{\textwidth}}{
\textbf{Resumen:} La rebeldía se configura cuando el demandado notificado no comparece ni contesta, y se declara a pedido del actor. Genera la continuación normal del proceso, la aplicación de las cargas del Art. 356 inc. 1°, una presunción de verdad de los hechos del actor y la posibilidad de trabar medidas cautelares, notificándose las resoluciones por ministerio de ley salvo la declaración de rebeldía y la sentencia, que se notifican por cédula. El rebelde que se presenta puede intervenir hacia adelante, salvo en los actos ya precluidos, y conserva la posibilidad de oponer la excepción de prescripción si justifica causa insuperable. Si la notificación de la demanda fue nula, el proceso se retrotrae y debe notificarse nuevamente.
} \\
\bottomrule
\end{tabularx}
\end{table}

% =====================================================
\chapter{Cuadro Cornell general de repaso}
% =====================================================

El siguiente cuadro reúne las preguntas y respuestas clave de toda la Unidad 10, tal como fueron sistematizadas en el material original, junto con una síntesis final de los contenidos estudiados.

\begin{longtable}{
    >{\raggedright\arraybackslash}p{0.32\textwidth}
    >{\raggedright\arraybackslash}p{0.58\textwidth}
}
\caption{Cuadro Cornell — Repaso general de la Unidad 10} \\
\toprule
\textbf{Preguntas / palabras clave} & \textbf{Notas / respuestas} \\
\midrule
\endfirsthead
\toprule
\textbf{Preguntas / palabras clave} & \textbf{Notas / respuestas} \\
\midrule
\endhead

\textbf{¿Qué es la contestación de demanda?} &
Es un acto procesal de postulación en el cual el demandado se defiende, pidiéndole a la jurisdicción que rechace la pretensión del actor invocada en su demanda. \\

\textbf{¿Cuál es el plazo para contestar en proceso ordinario?} &
15 días hábiles contados desde el día siguiente a la notificación. Existe un plazo de gracia de dos horas al inicio del día 16. Si el domicilio del demandado está a más de 100 km, el plazo se extiende un día por cada 200 km o fracción mayor a 100 km. \\

\textbf{Art. 356 inc. 1°: ¿qué debe hacer el demandado con los hechos?} &
Debe reconocer o negar categóricamente cada hecho invocado por el actor. El silencio o la negativa genérica facultan al juez a tener por admitidos esos hechos (es una facultad del juez). \\

\textbf{Art. 356 inc. 1°: ¿qué debe hacer el demandado con los documentos?} &
Debe admitir o negar categóricamente la autenticidad de cada documento. El silencio o la negativa genérica obligan al juez a tenerlos por reconocidos y auténticos (es una obligación del juez, no una facultad). \\

\textbf{¿Cómo se impugna un instrumento público?} &
Debe impugnarse al contestar, pero eso solo no basta. Para quitarle eficacia probatoria hay que promover un incidente de redargución de falsedad (Art. 395 CPCCN) dentro de los 15 días de haberlo impugnado. Sin el incidente, el instrumento público sigue haciendo plena fe. \\

\textbf{¿Qué es el allanamiento?} &
El acto procesal unilateral por el cual el demandado se somete a la pretensión del actor, sin necesitar la conformidad de éste. Puede ser total o parcial, expreso o tácito, y puede ejercerse en cualquier momento antes del dictado de la sentencia. \\

\textbf{¿Qué es la reconvención y cuál es su fundamento?} &
Es la demanda que el demandado dirige contra el actor dentro del mismo proceso, ejerciendo una pretensión autónoma. Su fundamento es la economía procesal: evitar dos procesos paralelos y el riesgo de sentencias contradictorias. Las pretensiones deben derivar de la misma relación jurídica o ser conexas. \\

\textbf{¿Qué es la rebeldía y cómo se declara?} &
Es la situación en que incurre el demandado debidamente notificado que no comparece y no contesta. Se declara a pedido del actor (no de oficio), una vez vencido el plazo para contestar. La resolución que la declara se notifica por cédula al domicilio real del demandado. \\

\textbf{¿Cuáles son los efectos de la rebeldía?} &
1) El proceso continúa normalmente (no se dicta sentencia automáticamente). 2) Se aplica el Art. 356 inc. 1°: el juez puede tener por admitidos los hechos y debe tener por reconocidos los documentos. 3) El estado de rebeldía implica presunción de verdad de los hechos del actor. 4) El actor puede solicitar medidas cautelares (Art. 63 y Art. 212 CPCCN). 5) Las resoluciones se notifican al rebelde por ministerio de ley. \\

\textbf{¿Qué puede hacer el demandado rebelde que se presenta tardíamente?} &
Puede ejercer todos los actos procesales a partir de su presentación y en las etapas hacia adelante (sistema de preclusión). No puede contestar demanda, ofrecer prueba ni oponer excepciones (esas etapas ya pasaron). Puede pedir la modificación, sustitución o reducción de medidas cautelares. Si demuestra que su rebeldía fue por causa insuperable (ej.: internación en coma), el juez puede levantar las medidas cautelares. \\

\textbf{¿Cuándo pueden incorporarse documentos nuevos?} &
En primera instancia: hasta que el expediente quede en ``autos para sentencia'' (Art. 335 CPCCN). En segunda instancia (Cámara): al fundar los recursos de apelación (Art. 260 CPCCN). Requisitos: que sean de fecha posterior a la demanda/contestación, o que, siendo anteriores, las partes juren no haber conocido su existencia. \\

\textbf{Principio de preclusión: ¿cómo rige en la contestación?} &
El sistema de preclusión implica que cada acto debe realizarse en la oportunidad prevista. Si no se oponen excepciones al contestar, ya no pueden oponerse. Si no se reconviene al contestar, ya no puede reconvenirse. Si no se recusa sin causa al juez al contestar, ya no puede hacerse. Las etapas no se reabren, salvo planteo de nulidad válido. \\

\midrule
\multicolumn{2}{p{\textwidth}}{
\textbf{Síntesis final:} En esta unidad se completó el estudio de la etapa introductoria del proceso judicial ordinario. Partiendo del principio de preclusión —que establece que cada acto procesal debe realizarse en su oportunidad bajo pena de perder el derecho a ejercerlo—, se analizó en detalle la contestación de demanda (acto de postulación por excelencia del demandado, regulada por el Art. 356 CPCCN), destacando las cargas de reconocer o negar categóricamente los hechos y los documentos, con consecuencias distintas para el silencio en uno y otro caso.

Se profundizó en las diferencias entre instrumentos públicos (que hacen plena fe y requieren un incidente de redargución de falsedad para ser impugnados) e instrumentos privados (que pierden eficacia probatoria con la mera negación categórica al contestar). Se estudió también el régimen de ofrecimiento de prueba en los escritos introductorios y las excepciones a esa regla: documentos nuevos (Art. 335), hechos no invocados (Art. 334) y hechos nuevos (Art. 365).

Se analizaron las tres actitudes especiales del demandado al contestar: el allanamiento (sometimiento unilateral a la pretensión del actor, total o parcial, expreso o tácito; modo anormal de terminar el proceso del Art. 307 CPCCN), la reconvención (demanda autónoma del demandado contra el actor dentro del mismo proceso, cuyo fundamento es la economía procesal, y que requiere la relación o conexidad entre las pretensiones) y la rebeldía (incumplimiento de las cargas de comparecer y contestar, que genera la presunción de verdad de los hechos del actor, habilita medidas cautelares al actor, y produce que las resoluciones se notifiquen por ministerio de ley, con excepción de la propia declaración de rebeldía y la sentencia, que se notifican por cédula al domicilio real).

Con el cierre de la etapa introductoria, el proceso queda listo para la audiencia preliminar del Art. 360 CPCCN, donde el juez identificará los hechos controvertidos y conducentes, y ordenará la producción de prueba para la etapa probatoria.
} \\
\bottomrule
\end{longtable}

\end{document}
